\subsection{Propositional Connectives: Truth Tables}

Truth-functional Connectives  (assume capital letters represent sentences)

Negation:
% \begin{displaymath}
\[
\begin{array}{|c|c|}
% |c c|c| means that there are three columns in the table and
% a vertical bar ’|’ will be printed on the left and right borders,
% and between the second and the third columns.
% The letter ’c’ means the value will be centered within the column,
% letter ’l’, left-aligned, and ’r’, right-aligned.
  A & \lnot A\\ % Use & to separate the columns
\hline % Put a horizontal line between the table header and the rest.
T & F\\
F & T\\
\end{array}
\]
% \end{displaymath}

Conjunction / And:

\[
\begin{array}{|c c|c|}
% |c c|c| means that there are three columns in the table and
% a vertical bar ’|’ will be printed on the left and right borders,
% and between the second and the third columns.
% The letter ’c’ means the value will be centered within the column,
% letter ’l’, left-aligned, and ’r’, right-aligned.
A & B & A \land B\\ % Use & to separate the columns
\hline % Put a horizontal line between the table header and the rest.
T & T & T\\
F & T & F\\
T & F & F\\
F & F & F\\
\end{array}
\]

Disjunction / Inclusive Or:

\[
\begin{array}{|c c|c|}
% |c c|c| means that there are three columns in the table and
% a vertical bar ’|’ will be printed on the left and right borders,
% and between the second and the third columns.
% The letter ’c’ means the value will be centered within the column,
% letter ’l’, left-aligned, and ’r’, right-aligned.
A & B & A \lor B\\ % Use & to separate the columns
\hline % Put a horizontal line between the table header and the rest.
T & T & T\\
F & T & T\\
T & F & T\\
F & F & F\\
\end{array}
\]

Conditional / If-then:

\[
\begin{array}{|c c|c|}
% letter ’l’, left-aligned, and ’r’, right-aligned.
A & B & A \imp B\\ % Use & to separate the columns
\hline % Put a horizontal line between the table header and the rest.
T & T & T\\
F & T & T\\
T & F & T\\
F & F & F\\
\end{array}
\]


Biconditional / If-then:

\[
\begin{array}{|c c|c|}
% letter ’l’, left-aligned, and ’r’, right-aligned.
A & B & A \bicon B\\ % Use & to separate the columns
\hline % Put a horizontal line between the table header and the rest.
T & T & T\\
F & T & F\\
T & F & F\\
F & F & T\\
\end{array}
\]

For Mendlelson's notation:

\begin{enumerate}
  \item Capital italic letters (sometimes w/ subscripts) are statement forms
  \item if $\B$ and $\C$ are statement forms, we can modify them with propositonal connectives to form other statements
\end{enumerate}

I don't feel the need to recreate the large truth tables, as they are clear enough in the text.

Note: If there are $n$ distinct letters in a statement form, there are $2^n$ possible assignments of truth values to the statement letters, meaning $2^n$ rows in the truth table.
